% ═══════════════════════════════════════════════════════════════════
%  Introduction générale
% ═══════════════════════════════════════════════════════════════════
\pageLiminaire{Introduction Générale}

La transformation numérique n'est plus l'apanage des grandes entreprises. Elle
touche aujourd'hui les commerces de proximité, les cabinets, les artisans et
les instituts, dont l'activité repose encore très souvent sur des outils
manuels : un carnet, un téléphone, une carte cartonnée. Ces outils ont une
qualité que l'on sous-estime --- ils sont immédiats, personne n'a besoin de les
apprendre --- et un défaut qui devient handicapant dès que l'activité croît :
ils ne conservent rien d'exploitable, ne se consultent qu'à un seul endroit, et
n'existent que pendant les heures d'ouverture.

Le secteur de la beauté et du bien-être illustre ce constat de façon
particulièrement nette. L'activité d'un institut s'organise autour d'une
ressource rare et non stockable --- un créneau de cabine --- que plusieurs
canaux se disputent en même temps : la cliente qui appelle, celle qui écrit sur
une messagerie instantanée, et celle qui se présente au comptoir. Le carnet
papier arbitre entre ces canaux, mais il ne peut le faire que si quelqu'un le
tient ouvert devant soi. Toute réservation prise ailleurs qu'à côté du carnet
est une réservation à risque.

C'est dans ce contexte que s'inscrit ce projet de fin d'études. \yani\ est un
centre de beauté installé à Casablanca, qui propose des prestations de soin et
vend des produits cosmétiques. Sa gérante y assure à la fois les soins, la
prise des rendez-vous, le suivi du stock et la fidélisation de sa clientèle.
L'objectif du projet est de porter l'ensemble de ces activités sur une
plate-forme numérique : une application mobile pour les clientes, un
back-office web pour l'institut, et une API qui porte les règles métier
communes aux deux.

Le travail présenté ici ne s'est pas limité à traduire un carnet papier en
formulaires. Une réservation en ligne pose des questions qu'un carnet ne pose
pas : que se passe-t-il lorsque deux clientes valident le même créneau à la
même seconde ? Comment un institut ouvert de 9\,h à 13\,h puis de 14\,h à
18\,h, à l'heure de Casablanca, expose-t-il ses disponibilités à une
application qui raisonne en temps universel ? Que devient l'historique
d'achat d'une cliente qui exerce son droit à l'effacement, alors même que cet
historique porte le chiffre d'affaires de l'institut ? Ces questions ont
structuré la conception autant que les fonctionnalités elles-mêmes, et elles
occupent une place proportionnée dans ce rapport.

Le présent document rend compte de ce travail. Il s'articule autour de quatre
chapitres.

Le \textbf{premier chapitre} présente l'organisme d'accueil, l'étude de
l'existant, la problématique qui en découle et la solution retenue, puis la
méthode de conduite de projet et sa planification.

Le \textbf{deuxième chapitre} est consacré à l'analyse et à la conception : la
spécification des besoins fonctionnels et non fonctionnels, les acteurs du
système, puis la modélisation UML --- diagrammes de cas d'utilisation,
diagrammes de séquence et diagramme de classes.

Le \textbf{troisième chapitre} expose l'étude technique : les langages, les
cadriciels et le système de gestion de base de données retenus, les
environnements de travail, et l'architecture technique de la solution, y
compris la chaîne de sécurité appliquée à chaque requête.

Le \textbf{quatrième chapitre} traite de la mise en \oe uvre : la structure
effective du code des trois applications, la stratégie de tests et
l'intégration continue, la présentation des interfaces réalisées, et enfin le
déploiement et l'exploitation de la plate-forme.

Une conclusion générale dresse le bilan du travail accompli, ses limites
assumées et les perspectives d'évolution de la solution.
